% !PRINTED 550
\weberpage{550}
% running head: 550 Siebzehnter Abschnitt. \S.\,151.
Da alle Zahlen aus $\mathfrak{a}\mathfrak{b}$ von der Form $\Sigma\alpha\beta$ sind, so ist zunächst
 klar, dass alle diese Zahlen durch $\varphi\psi$ theilbar sind.

Wenn wir aber nach \S.\,143, 4. das Functional $\varphi$ als grössten
gemeinschaftlichen Theiler zweier Zahlen $\alpha_1, \alpha_2$ darstellen, so
gehören diese Zahlen, als durch $\varphi$ theilbar, dem Ideal $\mathfrak{a}$ an,
und wir können, da es auf einen Einheitsfactor bei $\varphi$ nicht
ankommt,

\[
\varphi = \alpha_1 x_1 + \alpha_2 x_2
\]

setzen, wenn $x_1, x_2$ Variable sind. Ebenso können wir, wenn
$\beta_1, \beta_2$ zwei Zahlen aus $\mathfrak{b}$ und $y_1, y_2$ Variable bedeuten,

\[
\psi = \beta_1 y_1 + \beta_2 y_2
\]

setzen, und daraus ergiebt sich

\[
\varphi\psi = \alpha_1 \beta_1 x_1 y_1 + \alpha_1 \beta_2 x_1 y_2 + \alpha_2 \beta_1 x_2 y_1 + \alpha_2 \beta_2 x_2 y_2.
\]

Es ist also $\varphi\psi$ nach \S.\,143, 2. der grösste gemeinschaftliche
Theiler der vier Zahlen $\alpha_1\beta_1, \alpha_1\beta_2, \alpha_2\beta_1, \alpha_2\beta_2$, die dem Ideal
$\mathfrak{a}\mathfrak{b}$ angehören, und folglich ist $\varphi\psi$ der \gesperrt{grösste gemeinschaftliche
Theiler} aller Zahlen des Ideals $\mathfrak{a}\mathfrak{b}$.

Damit ist die gegenseitige Zuordnung der Functionale und
Ideale den Forderungen 1. bis 5. gemäss bewerkstelligt.

Den \gesperrt{Primfunctionalen} entsprechen bei dieser Zuordnung
\gesperrt{Primideale}, und die Zerlegung der Ideale in Primfactoren und
überhaupt die Gesetze der Theilbarkeit der Ideale ergeben sich
in völliger Uebereinstimmung mit den entsprechenden Sätzen
aus der Theorie der Functionale.

Bei dieser völligen Uebereinstimmung kann es zu keiner
Unzuträglichkeit führen, wenn wir das ganze System aller unter
einander associirten Functionale zu einem Gemeinbegriffe zusammenfassen
 und dafür den Namen \gesperrt{Ideal} brauchen.

Das System aller mit einer ganzen Zahl associirten Functionale
 ist dann ein \gesperrt{Hauptideal}.

Wenn irgend eine Zahl oder ein Functional durch die
Functionale eines Ideals theilbar ist, so nennen wir es \gesperrt{durch
das Ideal theilbar}, und wenn ein Functional $\varphi$ in Factoren
zerlegt ist, denen die Ideale $\mathfrak{a}, \mathfrak{b}, \ldots$ entsprechen, so setzen wir
auch, indem die Einheitsfactoren in der Bezeichnung weggelassen
 werden:

\weq{1}{\varphi = \mathfrak{a}\mathfrak{b} \ldots}

Eine Basis des Functionals ist zugleich eine Basis des Ideals,
und die absolute Norm des repräsentirenden Functionals stimmt
%CONT>>
